\documentclass[11pt,a4paper]{article}
\usepackage[margin=2.2cm]{geometry}
\usepackage[T1]{fontenc}
\usepackage{lmodern}
\usepackage{microtype}
\usepackage{amsmath,amssymb}
\usepackage{booktabs}
\usepackage{tabularx}
\usepackage{longtable}
\usepackage{array}
\usepackage{enumitem}
\usepackage{xcolor}
\usepackage[colorlinks=true,linkcolor=blue!50!black,urlcolor=blue!50!black]{hyperref}
\setlist{itemsep=2pt,topsep=3pt}
\setlength{\parskip}{4pt}
\setlength{\parindent}{0pt}
\renewcommand{\arraystretch}{1.18}
\newcolumntype{L}{>{\raggedright\arraybackslash}X}
\newcolumntype{P}[1]{>{\raggedright\arraybackslash}p{#1}}
\newcommand{\EST}{\textsuperscript{\textsf{est}}}
\newcommand{\keV}{\,\mathrm{keV}}
\newcommand{\GeV}{\,\mathrm{GeV}}
\newcommand{\TeV}{\,\mathrm{TeV}}
\newcommand{\abinv}{\,\mathrm{ab}^{-1}}
\newcommand{\fbinv}{\,\mathrm{fb}^{-1}}
\newcommand{\Bef}{B_{84}}
\newcommand{\must}{\textsf{\textbf{Must}}}
\newcommand{\should}{\textsf{Should}}
\newcommand{\could}{\textsf{Could}}

\title{\textbf{One publishable paper on electroweak quirks at the HL-LHC}\\[4pt]
\large A plan: what already stands, what to fix, and how to answer the referees}
\author{Planning document for the HL-LHC quirks study (paper v0.4)}
\date{12 November 2026}

\begin{document}
\maketitle

\begin{abstract}
\noindent
The study now ends in ``no working point, best expected limit 790 times the prediction, nothing excluded''. That result is
correct for its frozen rule, but it is not a publishable headline. The number measures the size of the simulated
background samples, not the physics of the search. This plan turns study and paper v0.4 into \emph{one} paper that covers
the whole work: the plane-thickness law, the detector-level search at pile-up 200, and a background \emph{estimate}
rather than a stack of zero-count bounds. Brute-force CPU does not solve the problem, since a hundred times more buys two
of the seven decades. The plan therefore spends effort on the bottlenecks that can actually be removed: the estimator,
the missing physics in the background model, the statistical treatment, signal statistics and acceptance, and the
paper's framing. Cloud computing is used only where it buys calendar time or disk, for an estimated total of a few
hundred US dollars. Everything stays open, with public tools and geometry and no CMS data.
\end{abstract}

\medskip
\noindent\textbf{Conventions.} Numbers without a mark are read from the study's tracked files, the paper v0.4 or the book
\emph{From Fields to Quirks}. A superscript \EST{} marks an estimate made for this plan, which nobody has run. Costs are
person-days (pd) and CPU-hours (CPU-h). Cloud prices assume 3 to 5 US cents per vCPU-hour, a round figure and not a
quotation. Priorities: \must{} (needed for a credible paper), \should{} (strongly improves it), \could{} (optional).

\tableofcontents
\newpage

%=====================================================================
\section{Where the study stands}

\subsection{The result in one paragraph}
A working point (WP) is the final selection used for the limit. The frozen rule takes the loosest of 23 cumulative cut
rungs (R01 to R23) whose 84th-percentile background $\Bef$ in $3\abinv$ is at most $0.11$ events, at that rung and at
every tighter one. A budget of $0.11$ means ``background-free'', and with $N_{\rm trig}=4.08\times10^7$ triggered events
it requires a fake rate near $2.7\times10^{-10}$ per trial. The lowest $\Bef$ on the ladder is $1.23\times10^6$ at R09,
$7.05$ decades too high, so no WP exists. With the pre-registered fallback selections, the best expected limit is 790
times the predicted cross section ($E$, $m_Q=200\GeV$, $\Lambda=1.5\keV$, efficiency $9/200$). Nothing is excluded, and the
existing bounds, all below about $180\GeV$ in this $\Lambda$ range, are not extended.

\subsection{Why that is not a physics statement}
\begin{itemize}
  \item From R08 on, every leading background term is a \textbf{zero-count bound}: no simulated event passes, so the term
        enters as $1.14/M$ with $M=500$ to $2000$, multiplied by the number of such events nature produces ($10^7$ to
        $10^9$). Such a bound depends only on the sample size, and tighter cuts earn nothing.
  \item The leading term at R09 is \textbf{decays in flight}, from a hand-made generator with estimated prong counts,
        because Fatras has no hadronic interactions and no decays in flight.
  \item The \textbf{pile-up estimate} by thinning failed its pre-registered closure test (ratio $1.50$ and $1.61$ at R04
        and R05, $z=+52.8$ and $+59.4$). It was patched by $\times2.41$, and beyond R06/R07 it rests on an untested
        factorisation, a factor-117 ambiguity between two ways of chaining the factors (``routes''), and an extrapolated
        trend width.
  \item As a result, \textbf{$\Bef$ rises past R09} as the cuts tighten, which no real background can do. Only the doubt
        grows.
\end{itemize}
So the study did not find a large background. It showed that counting simulated events cannot show the background is
small. That is the one fact the paper must fix or clearly frame.

\subsection{What brute force would cost}
Table~\ref{tab:brute} collects the study's own numbers. They are the reason this plan does not buy CPU for its own sake.

\begin{table}[h]
\centering\small
\caption{What raw statistics would cost (book, Ch.~27--28).}\label{tab:brute}
\begin{tabularx}{\linewidth}{@{}P{58mm}L@{}}
\toprule
step & effect \\
\midrule
Every background sample $\times100$ & $3.4\times10^4$ CPU-h, about \$1500\EST; two of the seven decades \\
Pile-up to budget by counting & about $10^9$ trials against $4.8\times10^4$; $3\times10^6$ CPU-h \\
Zero-count bound at budget, decays in flight & 35 times more events than the LHC will ever make of the process \\
Everything, if the zeros persist & $k_{\rm need}=1.1\times10^7$; about $4\times10^9$ CPU-h \\
\bottomrule
\end{tabularx}
\end{table}

\subsection{What a measured background would give}
The book's reach against an \emph{assumed} background (Ch.~29, \EST, best $\Lambda$ per mass, barrel only, $3\abinv$)
shows that the search does not need zero background. It needs a background that is \emph{estimated with a stated
width}:

\begin{table}[h]
\centering\small
\caption{Reach in GeV against an assumed background $B$ with log-normal width $\sigma_{\ln}$ (book Table 29.2, \EST).
Existing bounds in this range lie below about $180\GeV$; a dash means below $200\GeV$.}\label{tab:reach}
\begin{tabular}{@{}lrrrrrrrr@{}}
\toprule
$B$ & 0 & 1 & 10 & 30 & 100 & $10^3$ & $10^4$ & $10^5$\\
\midrule
$\sigma_{\ln}=0$   & 1026 & 971 & 866 & 799 & 717 & 548 & 393 & 243\\
$\sigma_{\ln}=0.2$ & 1026 & 970 & 854 & 760 & 622 & 325 & --  & -- \\
$\sigma_{\ln}=1$   & 1026 & 965 & 796 & 662 & 495 & --  & --  & -- \\
endcap, $g=1$      & 1265 & 1200 & 1062 & 979 & 897 & 729 & 554 & 396\\
\bottomrule
\end{tabular}
\end{table}

\textbf{The central target of this plan} follows from that table. Turn every leading zero-count bound into an estimate,
and bring the total background width to $\sigma_{\ln}\lesssim1$. A background of up to about $10^2$ events known to within
a factor of $e$ would then still give an expected exclusion well above the existing bounds (\EST).

%=====================================================================
\section{What is already publishable}
Table~\ref{tab:assets} lists the results that already stand, with the clean-up each needs. Together they are already
most of a paper. What is missing is a defensible background number.

\begin{longtable}{@{}P{34mm}P{61mm}P{59mm}@{}}
\caption{Results that already stand, with the clean-up each needs. N1 to N8 refer to the study's novelty audit
(\texttt{P2\_NOVELTY\_AUDIT.md}).}\label{tab:assets}\\
\toprule
result & what stands & clean-up before publication\\
\midrule\endfirsthead
\toprule
result & what stands & clean-up\\
\midrule\endhead
Plane-thickness law (P1, truth level) &
Analytic mechanism: the solenoid field acts as an electric field in the pair frame, the ends miss by $b_{\rm odd}$, and
the axis rotation carries $\ln(\ell_{\max}/b_{\rm odd})$. Law $\Delta s_B\propto\epsilon_B[\ln(1/\epsilon_B)+C]$, local
slope $c\simeq0.104$ at $1\keV$, $500\GeV$; blind numerics $0.105\pm0.002$; holds per event on $2.2\times10^5$ events;
multiple-scattering floor $10$--$28\,\mu$m, dominant above $1.2$--$1.9\keV$. It gives an analytic form of the rotation
that Li et al.\ found numerically. &
State the window as measured: there is a $+6.5\%$ departure inside the ``safe'' window at $5\keV$, and the law is followed
in 9 of 30 cells at 1--3\,keV, not ``at 1--3\,keV''. Report the failed $1.5\%$ population gate ($+2$--$3\%$) and the
failed straight-string ratio at $\gamma^*=2$ as such. Run the floor tests (under 1 CPU-h each) and the $B^*$-off test.
Fix the ``verified'' labels (P1 in \S\ref{sec:work}).\\
Hit-level chain at pile-up 200 (N1) &
First quirk ribbons through a full tracking-toolkit chain (ACTS/Fatras, digitisation, merging, CKF) of a public
Phase-2-like geometry at PU 200 normalised to TDR occupancies. Earlier studies used PU 50 or none. &
Detector-model robustness (B6). Linux reproducibility (C1). Public release of the configuration.\\
Signal efficiency and $\Lambda$ window (N6) &
Barrel $\varepsilon=0.085$--$0.31$ at the end of counting at $1$--$5\keV$; tracking loss factor $0.83$; $\Lambda$ window
lower edge $1.0$--$1.2\keV$ at PU 200, upper edge $4.7$--$6.6\keV$ without pile-up. &
More signal rows, the full chain at all 40 points (S1); the independent finder cross-check (C2); an endcap ladder (S3).\\
Background anatomy (N3) &
Random planes measured down 4.9 decades from R04 to R07 by direct count. First treatment of single-interaction
topologies (PV pairs, decays in flight, loopers, nuclear interactions, $V^0$) and 12 SM processes as quirk-plane
backgrounds. &
Replace bounds by estimates (B1--B4). Correct the attribution to Knapen et al., who name only beam-pipe conversions.\\
The simulation-statistics barrier (N4) &
Quantified: background-freedom needs $10^8$--$3\times10^9$ events per process, 2 to 35 times nature's own count. &
Keep it as a short section. It explains why the estimator in \S\ref{sec:work} is needed, and it is a useful warning for
every $B=0$ projection.\\
Timing category and ToF mass proxy (N7) &
Late-hit category; $M_{\rm est}$ median $2.26\,m_Q$ at PU 200, with 19\% below $0.5\,m_Q$ from in-time pile-up hits. &
Present it as a proxy, not a mass measurement.\\
Negative detector-level results (N8) &
The $\Lambda^{-2}$ separation handle is not measurable at PU 200, and the out-of-plane resolution ($23$--$30\,\mu$m)
equals the multiple-scattering floor. &
Keep them; they are honest and new.\\
Recast and reach &
At $B=0$: $1026\GeV$ in $3\abinv$. In $300\fbinv$: 3-event reach $665\GeV$ and 5-event reach $586\GeV$, against
$650\GeV$ for Knapen et al.\ in their set-up. &
Move it to the front, as validation that the signal side agrees with the literature.\\
Pre-registration (N5) &
Rule, ladder and fallback were frozen before any signal efficiency. &
Keep v1 as the frozen record; pre-register v2 before any rerun (\S\ref{sec:prereg}).\\
\bottomrule
\end{longtable}

%=====================================================================
\section{The one paper}\label{sec:paper}

\subsection{Message and title}
\textbf{Message.} \emph{A plane search for electroweak quirks in a Phase-2-like tracker at pile-up 200 works on the
signal side. With a background estimated, not bounded, it reaches $X$\,GeV, well beyond existing limits. We give the
thickness law that sets the ribbon width, the full background anatomy, and a data-driven recipe any experiment can apply.}

Working title: \emph{Electroweak quirks in a Phase-2 tracker at pile-up 200: plane thickness, detector-level search and
background-limited reach.}

Target: JHEP or Phys.\ Rev.\ D, about 35 to 45 pages with appendices, with a Zenodo record holding the code, configurations,
efficiency tables and likelihoods. SciPost Physics is a good open-access alternative.

\subsection{Structure}
\begin{enumerate}[label=\arabic*.]
  \item \textbf{Introduction.} The quirk signature, existing bounds, what is new, and the main result and figure, stated on
        page 1. Paper v0.4 buries its most meaningful sentence in \S10.
  \item \textbf{Model, samples and cross sections.} $E=(1,1)_{-1}$ and its scalar partner, $N_{\rm IC}=2$; NLO $K$-factors
        at 14\,TeV, and the $f_{200}$ uncertainty propagated (S2).
  \item \textbf{String dynamics and the thickness law (P1).} A compact main text; the derivation in an appendix.
  \item \textbf{Detector model.} Public geometry, pile-up overlay and top-up, timing layer, and robustness variations (B6).
  \item \textbf{Plane finder and selection.} Barrel ladder; endcap ladder if S3 is done; independent finder cross-check (C2).
  \item \textbf{Backgrounds.} Anatomy; why counting fails (N4, short); the estimators (B1--B4); closure tests; uncertainty
        model.
  \item \textbf{Results.} Efficiencies, the background ladder (v1 bounds against v2 estimates), expected limits on the
        $(m_Q,\Lambda)$ grid, and reach against background.
  \item \textbf{A data-driven background method for the experiments.} The road-rotation sideband, closed in simulation (B5).
  \item \textbf{Comparison} with Knapen et al., Li et al., FASER, D0 and the recasts.
  \item \textbf{Limitations,} stated plainly.
  \item \textbf{Conclusions.}
\end{enumerate}
Appendices: the thickness derivation; the pre-registration record (v1 frozen result, v2 changes and why); reproduction;
the limitation index; efficiency tables for recasting.

\subsection{Figures that must carry the message}
\begin{enumerate}
  \item An event display of a quirk ribbon at PU 200, labelled as the best of its sample.
  \item The thickness law with the multiple-scattering floor and the crossover.
  \item The efficiency map on the $(m_Q,\Lambda)$ grid, with the $\Lambda$ window.
  \item The background ladder: v1 zero-count bounds and v2 estimates with bands, against the budget.
  \item Closure plots for every estimator, at the rungs where direct counts exist.
  \item The expected limit against mass, with existing bounds overlaid.
  \item The reach against an assumed background (Table~\ref{tab:reach} as a curve), with the v2 estimate marked on it.
\end{enumerate}

\subsection{Writing clean-ups (W1)}
\begin{itemize}
  \item Remove internal labels that a referee cannot read: V4-3, C2-C, C4$'$, D13, A2FC, and \verb|\rep{...}| pointers to
        internal reports. Replace them with public supplementary material or plain words.
  \item Fix the known errors listed in \texttt{book/review/UNRESOLVED.md}:
    \begin{itemize}
      \item the $|\eta|<4$ cut stated for $f_{200}$, which the script does not apply;
      \item a 93\% vertex efficiency that the acceptance does not use;
      \item 0.966 quoted as the low end of the range, when it is the $0.45\keV$ entry (the low end at 0.6--5\,keV is 0.975);
      \item $0.105$ against $0.104$ called verified within 1.5\%, though the two come from different samples and on the
            same events the median is $+2.4\%$ high;
      \item $K_{\rm vac}$ listed among the blind numerics;
      \item ``linear in the endpoint proper time'', against the study's own audit;
      \item the attribution of a topology list to Knapen et al.
    \end{itemize}
  \item Tag every number as measured, estimated or unverified, as the book does. Keep negative results in the abstract.
  \item Disclose the AI-agent workflow, and the author's CMS membership together with the statement that only public
        information was used. Keep figures free of any CMS name or logo.
\end{itemize}

%=====================================================================
\section{Work plan: the bottlenecks that can be solved}\label{sec:work}

Each item states what it is, why it helps, its cost, and an \emph{acceptance criterion} fixed in advance in the v2
pre-registration. An item that misses its criterion is reported as failed and is not patched.

\subsection{Foundations}
\begin{description}[leftmargin=0pt,style=unboxed]
\item[C1 \must{} Restore and port.] Bring the 46\,GB of data back from the external drive and record where it lives.
  Port to Linux: add \texttt{linux-64} to the environments, remove macOS-only launchers, use a resumable runner with jobs
  under 30 minutes, and build a container. Reproduce the v1 numbers on Linux within statistics; results will not be
  bit-identical across compilers. \emph{Why:} without this no cloud or batch run is possible, and reproducibility is a
  referee question. \emph{Cost:} 3--5\,pd\EST. \emph{Criterion:} the v1 ladder reproduced within its statistical errors.
\item[C2 \must{} Finder cross-checks.] Run the planned independent brute-force search (V4-5) on 50 signal-plus-pile-up
  events and 50 pile-up trials, and complete the seed-threshold regression (V4-4). \emph{Why:} ``is your finder right?''
  is the first question a referee asks. \emph{Cost:} 2--3\,pd, under 50 CPU-h\EST.
\item[C3 \must{} Paper corrections.] The list in W1. \emph{Cost:} 2\,pd.
\end{description}

\subsection{Background: from bounds to estimates (the core)}
\begin{description}[leftmargin=0pt,style=unboxed]
\item[B1 \must{} Single-interaction topologies with real hadronic physics and a factorised estimator.]
  This is the leading term at the deep rungs.
  \begin{enumerate}
    \item \emph{Physics.} Use Geant4 for nuclear interactions in the material and for decays in flight ($\pi$, $K$ kinks,
          $V^0$), on a GDML or DD4hep description of the same cylinders, discs and material slabs as the ACTS geometry.
          Recommended hybrid: Geant4 produces the secondary vertices and their products, and Fatras propagates them, so
          the finder sees the same detector model as everywhere else. This replaces the hand-made generator with
          estimated prong counts.
    \item \emph{Estimator.} Write $P({\rm pass}\mid t)=P_{\rm geo}(t)\times C_{\rm PU}(t)$. $P_{\rm geo}$ is measured on large
          no-pile-up samples, which cost 1.5--12\,ms per event. $C_{\rm PU}$, how pile-up hits complete or spoil the
          candidate, is measured by overlaying those candidates on many pile-up events, at loose rungs where both
          factors have counts.
    \item \emph{Biased generation.} Only stiff configurations can pass a ribbon width of $4.7$--$9.4$\,mm, so generate
          the stiff tail with known weights and treat the soft bulk with an analytic sagitta bound that is validated
          where counts exist (\EST).
    \item \emph{Credit for purpose-built cuts.} See B3.
  \end{enumerate}
  \emph{Why:} it removes the decay-in-flight floor ($1.7\times10^5$ events from one bound in 500 events) and the
  referee's strongest objection, the missing hadronic physics. Gain 4--6 decades on these terms\EST{} if the zeros persist
  in the large samples. If they do not, the paper reports the first measured irreducible quirk-plane background, which is
  also a result. \emph{Cost:} 10--15\,pd; 1--5k CPU-h\EST, the one place where cloud is clearly useful.
  \emph{Criterion:} the factorised prediction closes against full pile-up samples at R05--R07 within a factor fixed in
  advance, and the Geant4 rates are cross-checked against the generator rates that v1 used.

\item[B2 \must{} Pile-up: repair the thinning closure and add an independent estimator.]
  Two independent routes to the same term:
  \begin{enumerate}
    \item \emph{Repair.} The study traced the closure failure to candidate merging, where candidates sharing at least half
          their hits count as one and merge too often at double density. Store the \emph{pre-merge} candidates at
          $\rho=2$ and re-apply merging per thinned subset, or define a density-independent candidate count. This needs
          about 0.4\,TB of storage, which the laptop (128\,GB free) could not hold: here money helps, through an external
          disk or cloud storage costing tens of dollars. It also supplies the R03 records that the R05 factorisation test
          needs and that do not exist.
    \item \emph{Independent estimator.} A semi-analytic occupancy model: per layer, the probability of at least two
          unassociated hits in a ribbon of the given width, measured to per-cent precision from existing events, combined
          with measured layer-pair correlations ($f_{\rm corr}=1.196$ exists) and closed against direct counts at
          R04--R07, a lever arm of 4.9 decades.
  \end{enumerate}
  \emph{Why:} it replaces the $\times2.41$ patch, the factor-117 route spread and the extrapolated trend width
  (2 to 6.3) with a tested estimate. \emph{Cost:} 10--15\,pd; about 500 CPU-h\EST.
  \emph{Criterion:} both estimators reproduce the direct counts at R04--R07 within the pre-set tolerance, and agree with
  each other at the deeper rungs within the quoted width.

\item[B3 \must{} Let the cuts act on what they were built for.] The rungs ``not a helix'' and ``not two straight lines''
  (R15--R20) reject exactly the prompt pairs and kinks that dominate the single-interaction terms, yet they earn no credit,
  because the one allowed bound per term is spent earlier. Add a displaced-vertex veto: a two-prong vertex at radius
  $>0$ removes conversions, $V^0$ and nuclear interactions. Measure each conditional on the large no-pile-up samples of B1,
  with its own closure test. \emph{Cost:} 3--5\,pd\EST.

\item[B4 \must{} Complete the Standard Model background.]
  Three steps:
  \begin{itemize}
    \item Produce hit-level samples for the eight processes now carried by proxies ($W(e\nu)$, $W(\mu\nu)$, dileptonic
          $t\bar t$, $t$-channel top, $tW$, $WW$, $WZ$, $ZZ$), which hold 15--45\% of the SM term.
    \item Add $Z(\tau\tau)$, lost-lepton $Z(\ell\ell)$ and multijets with fake missing energy at PU 200. They are absent
          now, and the multijets are the most exposed.
    \item Treat the hard-scatter tracks with the B1/B3 estimator instead of zero-count bounds.
  \end{itemize}
  \emph{Cost:} 5--8\,pd; about 300 CPU-h\EST.

\item[B5 \should{} The data-driven recipe, closed in simulation.] Rotate the road away from the missing-energy direction
  and keep every cut. Test, at R04--R07 on pile-up and SM samples with 20 roads per event, the two assumptions that have
  never been tested: the fake rate is flat in the angle to the missing energy, and roads in one event are independent.
  Quantify the part that is not flat, such as hadronic $\tau$ near the missing energy. \emph{Why:} the paper then gives
  the experiments a validated method, and the systematic floor of this method (the book's estimate is $\Bef\approx0.17$
  for an empty sideband) becomes the width used in the reach. \emph{Cost:} 5--8\,pd; about 500 CPU-h\EST.

\item[B6 \should{} Detector-model robustness.] Vary the hit resolution ($\times1.5$), the hit inefficiency, the material
  ($\times0.5$ and $\times2$; the material-variation sample takes 1.4\,h of wall clock) and the top-up limits, and
  propagate them to efficiency and background. \emph{Why:} it answers ``cylinders, not modules; Gaussian smearing;
  top-up $\times2.3$--$7.3$''. \emph{Cost:} 3--5\,pd; 100--200 CPU-h\EST.

\item[B7 \should{} Cosmics and beam halo.] Replace ``argued negligible'' with a number: public cosmic and halo rate
  measurements, combined with pointing to the primary vertex and timing. \emph{Cost:} 2\,pd\EST.
\end{description}

\subsection{Statistics}
\begin{description}[leftmargin=0pt,style=unboxed]
\item[T1 \must{} A proper likelihood.] Build a pyhf workspace (open) in which each background term is constrained by a
  Poisson auxiliary measurement of its simulated count. Zero counts are then handled by a Gamma constraint in the manner
  of Barlow and Beeston, not by $1.14/M$ entered as a log-normal. Correlate the shared nuisances across terms and
  categories, which v1 left uncorrelated. Remove the double counting in which the larger route is the central value and
  $\ln 117$ also enters the width. Replace the extrapolated trend width by the closure-derived systematics of B2. Publish
  the workspace. \emph{Why:} the sum of 21 separate 68\% bounds over-covers, and referees know it. \emph{Cost:} 5\,pd\EST.
\item[T2 \must{} A new selection rule.] Pre-register a choice among a few declared selections that optimises expected
  sensitivity, using an optimisation half and an evaluation half of the signal sample, so that the choice is not made on
  noise. Keep ``$\Bef\le0.11$'' as a reported \emph{background-free} flag, not as a requirement. \emph{Why:} the 0.11
  budget with $\sigma_{\ln}=2.3$ demands a median of 0.011 events. That is needed only to claim zero background, which
  nobody needs: $B=10$ costs 160\,GeV of reach (Table~\ref{tab:reach}). \emph{Cost:} 2\,pd.
\end{description}

\subsection{Signal}
\begin{description}[leftmargin=0pt,style=unboxed]
\item[S1 \must{} Signal statistics and one chain.] Ten times the rows at the 40 points, and the full chain at all of
  them. That retires the fast path, which may overestimate the efficiency by up to $1/0.87$. The best point now rests on
  9 events of 200. \emph{Cost:} about 400--600 CPU-h\EST, about \$20--30.
\item[S2 \must{} Theory inputs.] Compute NLO $K$-factors at 14\,TeV instead of transcribing them. Propagate the $f_{200}$
  uncertainty ($+12$--15\%, $-11$--13\%) into the limits. Re-shower with QED final-state radiation off the quirks, about
  15\,CPU-h, which was dropped only for budget. \emph{Cost:} 3\,pd.
\item[S3 \should{} Endcap finder and ladder.] 69\% of pairs with $|\eta_{\rm pair}|<2.5$ lie outside the barrel. Build
  endcap and mixed seeds and a separate endcap ladder with its own background from the B1/B2 estimators, frozen before its
  efficiency is seen. \emph{Gain:} signal $\times1.5$ to $3.2$\EST; reach at $B=10$ from 866 to 1062\,GeV for $g=1$.
  \emph{Caveats:} neither seed family has three layers at $1.0<|\eta|<1.3$, and the timing layer stops near
  $|\eta|\approx1.5$. \emph{Cost:} 15--25\,pd; 400--800 CPU-h\EST.
\item[S4 \could{} More efficiency.] A merging-tolerant rung for $\Lambda\gtrsim3\keV$, which was costed at about a week,
  and a lower cut on the seed's opening angle.
\item[S5 \could{} $m_Q$ and $\Lambda$ at PU 200.] A pile-up-robust separation estimator and rejection of in-time hits for
  the time of flight. Not required for the paper: state the negative result instead.
\end{description}

\subsection{P1 clean-ups}
\begin{description}[leftmargin=0pt,style=unboxed]
\item[P1 \must{}] Run the floor tests and the $B^*$-off test (under 1\,CPU-h each). Define the population estimator for
  $c$ in advance and report the failed gate as failed. State the validity window as measured. Keep the failed
  straight-string prediction visible. \emph{Cost:} 3--5\,pd.
\end{description}

%=====================================================================
\section{The v2 pre-registration}\label{sec:prereg}
Changing a rule after seeing a result looks like moving the goalposts, and referees will check for it. The defence is the
same discipline the study already practised.
\begin{enumerate}
  \item The v1 frozen result (no WP, 790$\times$) is kept and reported in an appendix with its hash, as the record of what
        pure counting gives.
  \item The v2 document is written and hashed \emph{before} any v2 sample is produced. It holds the estimators (B1--B4),
        their closure criteria, the likelihood (T1), the selection rule (T2), the endcap ladder (S3) if included, and what
        happens when a test fails.
  \item Every change is motivated by a background-method or physics argument stated in the document, never by a signal
        efficiency.
  \item Signal efficiencies stay blind until the background estimators have passed or failed their closure tests.
  \item Rows 1 to 4 of the book's ranking go into one freeze, because each freeze costs more than each change.
\end{enumerate}

%=====================================================================
\section{Compute: what to pay for and what not}
\begin{table}[h]
\centering\small
\caption{Compute for the plan (\EST). The laptop gives about 168 CPU-h per day at seven workers. A laptop CPU-h and a
cloud vCPU-hour are not the same unit, so read the totals as sizes.}\label{tab:cpu}
\begin{tabularx}{\linewidth}{@{}P{55mm}rL@{}}
\toprule
item & CPU-h & note\\
\midrule
B1 Geant4 vertices and no-PU tails & 1000--5000 & the main cloud item \\
B2 pile-up repair and occupancy & about 500 & plus 0.4--1\,TB of storage \\
B4 SM samples & about 300 & \\
B5 sideband, 20 roads & about 500 & \\
B6 robustness & 100--200 & \\
S1 signal $\times10$, full chain & 400--600 & \\
S3 endcap (if done) & 400--800 & \\
\midrule
total & about 3--8k & laptop: 3--7 weeks; 100 vCPU: 1.5--3.5 days; \$100--400 plus \$20--50 storage \\
\bottomrule
\end{tabularx}
\end{table}

\textbf{Worth paying for.} Cloud time for B1 and S1/S3, because it turns weeks of laptop wall clock into days. Disk for B2,
because the missing pre-merge records are a physics limitation that storage alone removes. A Linux container (C1) is the
prerequisite for both.

\textbf{Not worth paying for.}
\begin{itemize}
  \item Brute-force growth of the existing samples ($\times100$ for two decades).
  \item $10^9$ pile-up trials ($3\times10^6$ CPU-h).
  \item Rerunning the string at $\Lambda\ge15\keV$. It took 371 CPU-h, 64\% of the string's bill, and never reached the
        search; detector work needs 12--30 CPU-h of string.
  \item The $\rho=3$ sample: 128\,h for about 0.6 decade at R09, superseded by B2.
\end{itemize}

\textbf{Discarded by decision.} Pile-up scans (``background-free up to $\mu=X$'') and any use of CMS Open Data. Both are
outside the paper.

%=====================================================================
\section{Referee questions and the planned answers}
\small
\begin{longtable}{@{}P{4mm}P{54mm}P{36mm}P{56mm}@{}}
\toprule
\# & likely question & status in v0.4 & planned answer (item)\\
\midrule\endfirsthead
\toprule
\# & likely question & status in v0.4 & planned answer\\
\midrule\endhead
1 & Your background is a stack of zero-count bounds. What does the limit mean? & admitted; 790 is the reach of the set-up &
Estimators with closure tests; bounds kept only as the v1 record (B1--B4, T1).\\
2 & Fatras has no hadronic interactions; the leading background comes from a toy generator. & stated limitation &
Geant4 vertices fed into the same chain (B1).\\
3 & The thinning closure failed. Why trust any pile-up number? & patched $\times2.41$ &
Merging fixed and the closure re-tested; an independent occupancy estimator; both close at R04--R07 (B2).\\
4 & Factorisation across kinds of cut is untested. & degenerate test &
Conditionals measured on large samples, each with its own closure; correlations measured (B1--B3).\\
5 & Multijets with fake missing energy, $Z(\tau\tau)$ and lost-lepton $Z$ are missing. & not included & Added (B4).\\
6 & Eight SM processes are proxies. & 15--45\% of the SM term & Real samples (B4).\\
7 & The occupancy is topped up by $\times2.3$--$7.3$ over Fatras. & density systematic only & Robustness scan and propagation (B6).\\
8 & Cylinders, not modules; Gaussian smearing; no cluster merging. & stated & Resolution and inefficiency variations (B6); stated limits.\\
9 & Is the finder correct and near optimal? & cross-check not run & Brute-force comparison (C2).\\
10 & The best point rests on 9 events, and 31 points use a fast path. & stated & $\times10$ rows, full chain everywhere (S1).\\
11 & Why barrel only? & deadline decision & Endcap ladder (S3), or an explicit statement with the scaled gain if S3 is not done.\\
12 & You changed the rules after seeing the result. & v1 frozen & v2 hashed before any rerun; v1 in the appendix (\S\ref{sec:prereg}).\\
13 & Summed 68\% bounds over-cover; categories are uncorrelated. & simplified & pyhf likelihood with correlated nuisances and Gamma constraints (T1).\\
14 & Why a budget of 0.11? Nobody needs zero background. & frozen rule & Sensitivity-based rule; background-free reported as a flag (T2).\\
15 & Cosmics and beam halo? & argued, not simulated & Numerical estimate from public rates with pointing and timing (B7).\\
16 & Can the trigger fire on these events at Phase 2? & parametrised turn-on & Public L1/HLT turn-on curves, with a systematic (S2).\\
17 & $K$-factors and $f_{200}$ uncertainty? & transcribed; not propagated & NLO at 14\,TeV; propagated (S2).\\
18 & Classical Nambu--Goto string, $J=0$, no infracolour radiation? & stated assumption & Cite the validity regime; give the $\Lambda$ range where the assumptions hold.\\
19 & The thickness law's validity and its failed gate. & window overstated & Window as measured; failed gate and the failed ratio reported (P1).\\
20 & You cannot test the law at detector level. & stated (N8) & Keep it as a negative result: resolution equals the floor.\\
21 & Can you measure $m_Q$ or $\Lambda$? & no, at PU 200 & State the negative result; S5 optional.\\
22 & How does this compare with Knapen et al.\ and Li et al.? & recast agrees at $B=0$ & Comparison section: where the background-free assumption fails and what it costs.\\
23 & Is a data-driven estimate possible? & proposal only & Sideband closed in simulation (B5).\\
24 & Can I reproduce or recast this? & internal reports cited & Zenodo: code, container, configs, efficiency tables, pyhf workspace (C1, W1).\\
25 & Is this a CMS result? Was AI used? & labelling done & Disclosure paragraph; public information only (W1).\\
\bottomrule
\end{longtable}
\normalsize

%=====================================================================
\section{Outcomes and decision points}
\begin{table}[h]
\centering\small
\caption{Possible outcomes after v2 (\EST; reach from Table~\ref{tab:reach}, barrel, best $\Lambda$ per mass).}
\begin{tabularx}{\linewidth}{@{}P{30mm}P{62mm}L@{}}
\toprule
case & what happens & the paper\\
\midrule
Good & Leading terms become estimates of order 1--10 events, $\sigma_{\ln}\lesssim1$ & expected exclusion near 800\,GeV;
a strong pheno paper\\
Middle & Pile-up stays at $10^2$--$10^3$; single-interaction terms resolved & about 500\,GeV at $B=100$, $\sigma_{\ln}=1$:
still well above existing bounds\\
Real background & Large samples show passing events & the first measured irreducible quirk-plane background; the
reach at that $B$; the sideband recipe as the way out\\
Closure fails & An estimator misses its criterion & reported as failed; the paper becomes a method and anatomy paper
with P1 and the efficiencies, which is still publishable but less strong\\
\bottomrule
\end{tabularx}
\end{table}

\textbf{Decision points.}
\begin{enumerate}
  \item After C1 and C2: is the v1 chain reproducible on Linux, and does the finder pass its cross-check?
  \item After the B1 and B2 development, before the freeze: do the estimators close on existing samples at loose rungs? If
        not, redesign before freezing rather than after.
  \item After v2 production: unblind the background, then the signal.
  \item Include S3 (endcap) only if B1 and B2 are on schedule. Otherwise state the barrel restriction and its scaled gain.
\end{enumerate}

%=====================================================================
\section{Timeline}
\begin{table}[h]
\centering\small
\caption{Timeline (\EST). Calendar weeks assume one person, cloud for the heavy runs and laptop for the rest.}
\begin{tabularx}{\linewidth}{@{}P{34mm}P{20mm}L@{}}
\toprule
phase & weeks & content\\
\midrule
0 Foundations & 1--2 & C1, C2, C3; v1 reproduced on Linux\\
1 Method development & 3--4 & B1--B3 designs, closure on existing samples; T1; P1 tests; endcap prototype; signal blind\\
2 Freeze & 1 & v2 pre-registration written, reviewed, hashed\\
3 Production & 1--2 & B1, B2, B4, B5, B6, S1 (and S3) on the cloud, about \$100--400\\
4 Unblinding & 1 & background closure verdicts, then the signal; limits and reach\\
5 Paper & 2--3 & rewrite to the structure of \S\ref{sec:paper}; adversarial internal review against the referee table; Zenodo release\\
\midrule
total & 9--13 & core (without S3): about 8--10 weeks\\
\bottomrule
\end{tabularx}
\end{table}

\section*{Sources}
Study paper v0.4 (\texttt{stage2\_C001\_quirks/theory\_track/paper/sections/}, chiefly \S\S0, 10, 12, 13 and App.~C);
novelty audit \texttt{P2\_NOVELTY\_AUDIT.md}; book chapters 25--29 (\texttt{book/chapters/}); the paper corrections in
\texttt{book/review/UNRESOLVED.md}; a brainstorm by a second model, whose suggestions are marked \EST{} here. Nothing was
rerun for this plan.

\end{document}
